\section{Alternating characters and quasi-divisorial valuations}
\label{sec:local}

We first recover valuations from the degree-two relations.  The local
theory is most naturally expressed on the dual of the multiplicative
group.  Its multiplicative formulation also applies to residue fields
of characteristic~\(\ell\), which will occur in the argument.

For any field~\(F\), put
\[
 A_F=\Hom(F^\times,\Lambda)
     =\Hom_\Lambda(V_F,\Lambda),\qquad \Lambda=\F_\ell.
\]
Here the definition of~\(V_F\) is used in every characteristic.  Give
\(A_F\) the topology of pointwise convergence, with~\(\Lambda\) discrete.
Choosing a basis of~\(V_F\) identifies~\(A_F\) with a product of copies
of~\(\Lambda\).  In particular it is compact, and evaluation identifies
\(V_F\) with~\(\Hom_{\mathrm{cont}}(A_F,\Lambda)\): a continuous linear
form on that product depends on only finitely many coordinates.

A pair~\(f,g\in A_F\) is \emph{alternating} if
\begin{equation}
 f(x)g(1-x)=f(1-x)g(x)\qquad(x\in F\setminus\{0,1\}).
 \label{eq:local-alternating}
\end{equation}
A subset is alternating if every pair of its elements is alternating.
The determinant functional
\[
 [x]\otimes[y]\longmapsto f(x)g(y)-f(y)g(x)
\]
annihilates~\(R_F\) exactly when~\(f,g\) are alternating.  Thus the
isomorphism~\(\Theta\) in Theorem~\ref{thm:main} induces a continuous
linear isomorphism
\[
 \Phi=\Theta^*:A_L\xrightarrow{\ \sim\ }A_K
\]
that preserves alternating pairs in both directions.

We consider valuations up to equivalence and write~\(v\leq w\) when
\(v\) is a coarsening of~\(w\).  Write~\(U_v\) and~\(U_v^1\) for its
units and principal units, \(Fv\) for its residue field, and~\(vF\) for
its value group.  The \emph{minimized inertia and decomposition spaces}
are
\[
 I_v=U_v^\perp\subset D_v=(U_v^1)^\perp\subset A_F,
\]
where perpendiculars refer to evaluation.  A subset of~\(A_F\) is
\emph{valuative} if it is contained in some~\(I_v\).  Restriction to
units gives a continuous map
\[
 \rho_v:D_v\longrightarrow A_{Fv}
\]
with kernel~\(I_v\).  Surjectivity will be established for the valuations
we use below.

\subsection{The established local input}

The following form of alternating-pair valuation theory collects
\cite[Fact~2.1, Lemma~2.5, Theorem~2.6, and Corollary~2.7]{Topaz2020}.
We state it for finite coefficients; the field in this statement need
not have characteristic different from~\(\ell\).

\begin{theorem}[Alternating-pair valuation theory]
\label{thm:local-input}
Let~\(F\) be a field and set
\(A_F^{\pm}=\{f\in A_F:f(-1)=0\}\).
\begin{enumerate}
\item Every valuative subset~\(\Sigma\subset A_F\) has a unique coarsest
valuation~\(v_\Sigma\) with~\(\Sigma\subset I_{v_\Sigma}\).  In particular,
\(v_\Sigma\) coarsens every valuation whose inertia contains~\(\Sigma\).
For such a valuation~\(w\), this coarsening is obtained by quotienting
\(wF\) by its largest convex subgroup contained in
\(w(\bigcap_{f\in\Sigma}\ker f)\).
\item If~\(S\subset A_F^{\pm}\) is an alternating linear subspace, its
valuative elements form a valuative subspace~\(S_0\) of codimension at
most one.  Moreover~\(S\subset D_{v_{S_0}}\).
\item If~\(\Sigma\subset A_F^{\pm}\) is valuative and~\(g\in A_F^{\pm}\)
is alternating with every element of~\(\Sigma\), then
\(g\in D_{v_\Sigma}\).
\end{enumerate}
\end{theorem}

The coefficient convention in \cite[Section~2.1]{Topaz2020} explicitly
allows~\(\Z/\ell\); its fraction field is then~\(\F_\ell\), so its
character space is exactly~\(A_F\).  The case~\(\ell=2\) is included
in the local theorem.  For the fields over algebraically closed
constants considered here, every character kills the constants and
hence~\(-1\), so~\(A_F^{\pm}=A_F\).

\subsection{Value groups and residue characters}

For the remainder of this section, let~\(F/\kappa\) be a finitely
generated extension of an algebraically closed field, and put
\(s=\trdeg(F/\kappa)\).  No restriction is imposed on its characteristic.
We first record the valuation facts that turn the local theorem into
an intrinsic description of codimension-one valuations.

\begin{lemma}
\label{lem:local-abhyankar}
For every valuation~\(v\) of~\(F\), the field~\(\kappa v\) is
algebraically closed, the group~\(vF/v\kappa\) is torsion free, and
\begin{equation}
 \rrank(vF/v\kappa)+\trdeg(Fv/\kappa v)\leq s.
 \label{eq:local-abhyankar}
\end{equation}
Here~\(\rrank(G)=\dim_\Q(G\otimes_\Z\Q)\).  If equality holds, then
\(Fv/\kappa v\) is finitely generated and~\(vF/v\kappa\) is a finitely
generated free abelian group.  Equality passes to every coarsening
and to the induced residue valuation.  Conversely, equality for a
valuation and an induced residue valuation gives equality for their
composition.
\end{lemma}

\begin{proof}
The value group of an algebraically closed field is divisible, and its
residue field is algebraically closed.  If~\(n\gamma\in v\kappa\) for
\(\gamma\in vF\), choose~\(\delta\in v\kappa\) with
\(n\delta=n\gamma\).  Ordered abelian groups are torsion free, so
\(\gamma=\delta\).  This proves torsion-freeness of the relative group.

Choose elements~\(x_1,\ldots,x_a\in F^\times\) whose values are
rationally independent modulo~\(v\kappa\), and units
\(y_1,\ldots,y_b\) whose residues are algebraically independent
over~\(\kappa v\).  These elements are algebraically independent
over~\(\kappa\).  Indeed, in a putative polynomial relation, distinct
monomials in the~\(x_i\) have distinct values modulo~\(v\kappa\).
Within a coefficient polynomial in the~\(y_j\), its terms of least
value cannot cancel, by algebraic independence of their residues after
rescaling by a constant.  The nonzero summands obtained by grouping
according to monomials in the~\(x_i\) have distinct values, so their
sum cannot vanish.  This proves
\eqref{eq:local-abhyankar} by taking maximal such families.

Suppose equality holds and choose these families with~\(a+b=s\).
Then~\(F\) is finite over
\(F'=\kappa(x_1,\ldots,x_a,y_1,\ldots,y_b)\).  The same least-value
calculation gives
\[
 vF'/v\kappa\cong\Z^a,\qquad
 F'v=\kappa v(\bar y_1,\ldots,\bar y_b).
\]
The fundamental inequality for the finite extension~\(F/F'\) shows
that~\([vF:vF']\) and~\([Fv:F'v]\) are finite.  Thus the residue
extension is finitely generated, while the relative value group is
finitely generated and torsion free, hence free.

For the assertion about composition, write~\(v=w\circ u\), meaning
that~\(u\) is a coarsening and~\(w\) is the induced valuation on~\(Fu\).
The value groups fit into the exact sequence
\[
 0\longrightarrow w(Fu)/w(\kappa u)
 \longrightarrow vF/v\kappa
 \longrightarrow uF/u\kappa\longrightarrow0.
\]
Apply~\eqref{eq:local-abhyankar} first to~\(u\), and then to~\(w\)
on~\(Fu/\kappa u\).  The sum of the two inequalities is precisely the
inequality for~\(v\).  Consequently equality for~\(v\) forces equality
in both, and equality in both forces equality for~\(v\).  This use of
the inequality on~\(Fu/\kappa u\) is legitimate even before finite
generation is known: the algebraic-independence argument proving it
applies to arbitrary field extensions of finite transcendence degree.
\end{proof}

\begin{definition}
A \emph{quasi-divisorial valuation}, or \emph{quasi-prime divisor}, of
\(F/\kappa\) is a valuation minimal under coarsening among those with
\[
 vF/v\kappa\cong\Z,\qquad \trdeg(Fv/\kappa v)=s-1.
\]
A quasi-prime divisor trivial on~\(\kappa\) is called a \emph{prime
divisor}, or a \emph{true prime divisor}.
A \emph{quasi-prime \(r\)-divisor} is a composition of~\(r\)
successive quasi-prime divisors on the resulting residue fields.
\end{definition}

Lemma~\ref{lem:local-abhyankar} shows that for a quasi-prime
\(r\)-divisor~\(v\), the residue field is a function field of
transcendence degree~\(s-r\) over~\(\kappa v\), and
\begin{equation}
 vF/v\kappa\cong\Z^r.
 \label{eq:local-flag-rank}
\end{equation}
The isomorphism concerns abstract abelian groups; no ordering of
\(\Z^r\) is specified.

\begin{lemma}
\label{lem:residue-duality}
Suppose~\(vF/v\kappa\) is free abelian of finite rank.  Then restriction
to units induces a canonical topological isomorphism
\begin{equation}
 D_v/I_v\xrightarrow{\ \sim\ }A_{Fv}.
 \label{eq:local-residue-duality}
\end{equation}
Two elements of~\(D_v\) are alternating if and only if their images
in~\(A_{Fv}\) are alternating.  If~\(w\) is a valuation of~\(Fv\), then
\begin{equation}
 I_{w\circ v}=\rho_v^{-1}(I_w),\qquad
 D_{w\circ v}=\rho_v^{-1}(D_w),
 \label{eq:local-composition}
\end{equation}
where the inverse images are taken inside~\(D_v\).
\end{lemma}

\begin{proof}
There is an exact sequence of abelian groups
\[
 1\longrightarrow (Fv)^\times/(\kappa v)^\times
 \longrightarrow F^\times/(\kappa^\times U_v^1)
 \longrightarrow vF/v\kappa\longrightarrow0.
\]
For the kernel assertion, rescale by a constant an element whose value
lies in~\(v\kappa\), and then take its residue.  The last group is free,
so the sequence splits as a sequence of abstract groups.  Dualizing
into~\(\Lambda\) is consequently exact.  Characters kill the
algebraically closed constants on both sides, giving the asserted
surjection~\(\rho_v\) with kernel~\(I_v\).  Its induced bijection is a
homeomorphism because its source is compact and its target Hausdorff.

For alternation, if~\(v(x)>0\), then~\(1-x\in U_v^1\), and the
identity~\eqref{eq:local-alternating} holds for any pair in~\(D_v\).
The case~\(v(1-x)>0\) is the same.  If~\(v(x)<0\), then
\((1-x)/(-x)\in U_v^1\), so every character in~\(D_v\) takes the same
value on~\(x\) and~\(1-x\).  The remaining case is
\(v(x)=v(1-x)=0\), when the identity is exactly its residue-field
counterpart.  Conversely, lift any residue element different from
zero and one to test the residue identity.

Finally~\(U_v^1\) is contained in both~\(U_{w\circ v}\) and
\(U_{w\circ v}^1\).  Under reduction of~\(v\)-units, these two groups
map respectively onto~\(U_w\) and~\(U_w^1\).  Their annihilators
therefore give~\eqref{eq:local-composition}.
\end{proof}

We will also use the following consequence of valuation approximation.
It explains why intersections of decomposition spaces yield inertia.

\begin{lemma}
\label{lem:local-approximation}
Let~\(v_1,v_2\) be incomparable valuations of a field, and let~\(u\)
be their finest common coarsening.  Then
\[
 U_{v_1}^1U_{v_2}^1=U_u,
 \qquad D_{v_1}\cap D_{v_2}=I_u.
\]
\end{lemma}

\begin{proof}
The induced valuations~\(w_1,w_2\) on the residue field~\(Fu\) are
independent.  For~\(a\in(Fu)^\times\), the approximation theorem gives
\(b\in(Fu)^\times\) with
\[
 w_1(b-1)>0,\qquad w_2(b-a)>w_2(a).
\]
Thus~\(b\in U_{w_1}^1\) and~\(a/b\in U_{w_2}^1\), proving
\(U_{w_1}^1U_{w_2}^1=(Fu)^\times\).  Reduction maps~\(U_{v_i}^1\)
onto~\(U_{w_i}^1\), and both groups~\(U_{v_i}^1\) contain~\(U_u^1\).
Lifting the product identity gives the first equality; taking
annihilators gives the second.
\end{proof}

\subsection{Recovering dimension and quasi-divisors}

The next argument follows the alternating-space method of
\cite[Fact~3.2 and Theorem~3.3]{Topaz2020}.  We give the details so that
the argument remains available in characteristic~\(\ell\).

\begin{proposition}
\label{prop:local-dimension}
The maximum dimension of an alternating subspace of~\(A_F\) is~\(s\).
If~\(S\subset A_F\) is alternating of dimension~\(s\), and~\(v\) is
the valuation associated to its valuative part~\(S_0\), then
\(S_0=I_v\) and equality holds in~\eqref{eq:local-abhyankar} for~\(v\).
\end{proposition}

\begin{proof}
For any torsion-free abelian group~\(G\) of finite rational rank,
\[
 \dim_\Lambda(G/\ell G)\leq\rrank(G).
\]
Indeed, representatives of linearly independent classes modulo~\(\ell\)
are rationally independent: an integral relation can be divided by
the greatest common divisor of its coefficients, using torsion-freeness,
and the resulting relation has a coefficient nonzero modulo~\(\ell\).
Since characters of~\(vF\) kill its divisible subgroup~\(v\kappa\),
this proves
\[
 \dim I_v\leq\rrank(vF/v\kappa).
\]

Now let~\(S\) be alternating.  By Theorem~\ref{thm:local-input}, its
valuative part~\(S_0\) is a subspace of codimension at most one and,
for its associated valuation~\(v\),
\(S_0\subset I_v\) and~\(S\subset D_v\).  Moreover
\(S\cap I_v=S_0\), since every element of~\(I_v\) is valuative.
If~\(S=S_0\), then
\[
 \dim S\leq\dim I_v\leq\rrank(vF/v\kappa)\leq s.
\]
If~\(S\ne S_0\), an element of~\(S\setminus S_0\) has nonzero
image under~\(\rho_v\), whose kernel is always~\(I_v\).
Thus~\(A_{Fv}\ne0\).  Since~\(\kappa v\) is algebraically closed,
this implies~\(\trdeg(Fv/\kappa v)\geq1\), and hence
\[
 \dim S=\dim S_0+1
 \leq\rrank(vF/v\kappa)+\trdeg(Fv/\kappa v)\leq s.
\]
If~\(\dim S=s\), equality throughout the applicable chain gives
\(S_0=I_v\) and equality in~\eqref{eq:local-abhyankar}.

To attain the bound, choose a transcendence basis of~\(F/\kappa\),
take a full coordinate flag on its rational function field, and prolong
the associated valuation to~\(F\).  It is trivial on~\(\kappa\), and
its value group is a finite-index extension of~\(\Z^s\), hence is
abstractly~\(\Z^s\).  Its inertia has dimension~\(s\) and is
alternating by Lemma~\ref{lem:residue-duality}.  For~\(s=0\), the
field is~\(\kappa\), its character space is zero, and the assertion
has the same interpretation.
\end{proof}

We can now identify the one-dimensional inertia spaces solely through
maximal alternating subspaces.  Once inertia is identified, alternation
with it identifies decomposition as well.

\begin{proposition}
\label{prop:quasi-recovery}
Assume~\(s\geq2\).  A line~\(H\subset A_F\) is the inertia of a
quasi-divisorial valuation if and only if
\begin{equation}
 H=S_1\cap S_2
 \quad\text{for alternating subspaces }S_1,S_2\subset A_F
 \text{ of dimension }s.
 \label{eq:local-inertia-intersection}
\end{equation}
The quasi-divisorial valuation~\(v\) with~\(I_v=H\) is unique.  For any nonzero
\(h\in H\), its decomposition space is
\begin{equation}
 D_v=\{g\in A_F:(h,g)\text{ is alternating}\}.
 \label{eq:local-decomposition-recognition}
\end{equation}
\end{proposition}

\begin{proof}
Suppose first that~\(v\) is quasi-divisorial.  On~\(Fv/\kappa v\)
choose two independent full discrete flag valuations~\(w_1,w_2\),
trivial on~\(\kappa v\).  To obtain them, start with coordinate flags
on a rational subfield of transcendence degree~\(s-1\), with distinct
first prime divisors, and prolong to the finite extension~\(Fv\).
Independence survives prolongation.  Indeed, a valuation on a finite
field extension that restricts trivially to the smaller field is
trivial, since its value group is then finite and hence zero.  Thus a
common nontrivial coarsening of~\(w_1,w_2\) would restrict to a common
nontrivial coarsening of the original independent flag valuations,
which is impossible.

The inertia spaces~\(I_{w_1},I_{w_2}\) have zero intersection by
approximation.  Set~\(v_i=w_i\circ v\) and~\(S_i=I_{v_i}\).
Each~\(v_i\) has relative value group~\(\Z^s\), so its inertia is
alternating of dimension~\(s\).  By~\eqref{eq:local-composition}, the
intersection~\(S_1\cap S_2\) is~\(I_v\).

Conversely, suppose~\eqref{eq:local-inertia-intersection} holds.
Let~\(S_{i,0}\) be the valuative part of~\(S_i\), with associated
valuation~\(v_i\).  Proposition~\ref{prop:local-dimension} gives
\[
 S_{i,0}=I_{v_i},\qquad S_i\subset D_{v_i},
 \qquad \dim I_{v_i}\geq s-1\geq1,
\]
and~\(v_i\) has equality in~\eqref{eq:local-abhyankar}.
We first show that~\(H\) is valuative.  If, for example,
\(v_1\leq v_2\), then
\[
 0\ne I_{v_1}\subset S_1\cap S_2=H,
\]
so~\(H=I_{v_1}\).  The other comparable case is symmetric.  If the
valuations are incomparable, Lemma~\ref{lem:local-approximation}
puts~\(H\) inside the inertia of their finest common coarsening.
In either case every element of~\(H\) is valuative.  Because it lies
in each~\(S_i\), it then lies in each~\(S_{i,0}\), giving
\[
 H=I_{v_1}\cap I_{v_2}.
\]

Let~\(u=v_H\).  By Theorem~\ref{thm:local-input}, it coarsens each
\(v_i\), and therefore
\[
 H\subset I_u\subset I_{v_1}\cap I_{v_2}=H.
\]
Equality in~\eqref{eq:local-abhyankar} passes to~\(u\).
Its relative value group is free, so~\(\dim I_u=1\) implies
\(uF/u\kappa\cong\Z\) and~\(\trdeg(Fu/\kappa u)=s-1\).
If a coarsening~\(u'\leq u\) still has relative rank one, its
one-dimensional inertia is contained in~\(I_u\), hence equals~\(H\).
The defining coarseness of~\(v_H=u\) then gives~\(u\leq u'\), so
\(u=u'\).  Thus~\(u\) is quasi-divisorial.

For uniqueness, if~\(v\) is any quasi-divisor with~\(I_v=H\), its
associated coarsening~\(v_H\) has the same inertia.  Equality in
\eqref{eq:local-abhyankar} passes to that coarsening, so it too has
relative rank one and residue transcendence degree~\(s-1\).
Minimality of~\(v\) forces~\(v=v_H\).

Finally, an element of~\(D_v\) is alternating with~\(h\) by
Lemma~\ref{lem:residue-duality}, since~\(h\) has zero residue image.
Conversely, an element alternating with~\(h\) belongs to
\(D_{v_H}=D_v\) by Theorem~\ref{thm:local-input}.
\end{proof}

\begin{lemma}
\label{lem:local-flag-uniqueness}
If~\(v\) is a quasi-prime~\(r\)-divisor, with~\(r\geq1\), then its
value group has no nonzero~\(\ell\)-divisible convex subgroup.
Moreover~\(v=v_{I_v}\); in particular its inertia determines the
valuation.
\end{lemma}

\begin{proof}
First suppose~\(r=1\).  An~\(\ell\)-divisible convex subgroup~\(C\)
of~\(vF\) maps to zero in~\(vF/v\kappa\cong\Z\), and hence lies
in~\(v\kappa\).  Coarsening by~\(C\) therefore leaves the relative
value group unchanged.  Lemma~\ref{lem:local-abhyankar} shows that the
residue transcendence degree remains~\(s-1\), so minimality
of~\(v\) forces~\(C=0\).

For a composite flag, the value group of its final quasi-prime step
is a nonzero convex subgroup~\(B\subset vF\) with the property just
proved.  If~\(C\subset vF\) were a nonzero~\(\ell\)-divisible convex
subgroup, then~\(B\cap C\ne0\), since convex subgroups are linearly
ordered by inclusion.  This intersection is~\(\ell\)-divisible:
division by~\(\ell\) in~\(C\) stays in~\(B\) by convexity.  It is
also convex in~\(B\), a contradiction.

Put~\(\Gamma=vF\).  Since~\(\Gamma/v\kappa\cong\Z^r\) and
\(v\kappa\) is divisible, the common kernel in~\(\Gamma\) of all
characters in~\(I_v\) is
\[
 v\kappa+\ell\Gamma=\ell\Gamma.
\]
Every convex subgroup contained in~\(\ell\Gamma\) is itself
\(\ell\)-divisible: an~\(\ell\)-th part exists in~\(\Gamma\), and
convexity places it in the subgroup.  The first assertion therefore
shows that the largest such convex subgroup is zero.  The
coarsening description in Theorem~\ref{thm:local-input} now gives
\(v=v_{I_v}\).
\end{proof}

\begin{corollary}
\label{cor:flag-matching}
For the fields~\(K/k,L/l\) of Theorem~\ref{thm:main}, a compatible
isomorphism~\(\Theta\) forces
\(\trdeg(K/k)=\trdeg(L/l)=d\), say.  Its dual~\(\Phi\) matches the
inertia--decomposition pairs of quasi-prime divisors, and, successively,
the pairs of quasi-prime~\(r\)-divisors for~\(1\leq r\leq d-1\).
For a matched pair of valuations~\(v\) on~\(K\) and~\(w\) on~\(L\),
it induces a topological linear isomorphism
\[
 A_{Lw}\xrightarrow{\ \sim\ }A_{Kv}
\]
preserving alternating pairs in both directions.
\end{corollary}

\begin{proof}
The maximum alternating dimension is preserved by~\(\Phi\), giving
equality of transcendence degrees by
Proposition~\ref{prop:local-dimension}.  Proposition~\ref{prop:quasi-recovery}
then identifies corresponding quasi-prime inertia and decomposition
spaces.  Quotienting by inertia and using
Lemma~\ref{lem:residue-duality} gives the indicated isomorphism on
residue characters, with its alternating relation.

As long as the residual transcendence degree is at least two, apply
Proposition~\ref{prop:quasi-recovery} again there.  The inverse-image
identities~\eqref{eq:local-composition} identify the resulting pairs
with those of compositions on the original fields.
Lemma~\ref{lem:local-flag-uniqueness} shows that each such inertia
determines its composite valuation, so this matching is independent
of any choice of a presentation as a flag.  At each step
Lemma~\ref{lem:local-abhyankar} supplies finitely generated residue
fields over algebraically closed residue constants.  The local
results apply even if one of these fields has characteristic~\(\ell\).
Iteration ends after~\(d-1\) steps, with function fields of curves.
\end{proof}
